\section{Open Work}

The current code checks structured diffs. It parses JSON, validates hunks,
checks the base file, applies diffs, detects conflicts, models linear histories,
and checks merges. The following work is still open.

\subsection{Monadic Pipeline}

Add a monadic version of the pipeline. It must run parsing, schema checks, base
checks, conflict checks, and merge checks in one sequence. The current code
exposes these steps as separate functions.

\subsection{Invalid Input Flags}

Add a result type for bad input from an agent. It must separate JSON parse
errors, schema errors, and base-alignment errors. The current code returns
\code{Except String}; callers still have to classify the failure.

\subsection{Merge Conflict Flags}

Add a result type for merge conflicts. The current code computes conflict
witnesses with \code{conflictWitnesses} and returns them through
\code{MergeResult.conflict}. Callers still have to present that result as a
clear merge-conflict flag.

\subsection{Agent Runner}

There is no file that runs agents. A separate runner can call
\code{parseDiff}, \code{Diff.checkSchema}, \code{Diff.checkAgainstBase},
\code{nWayCompatible}, and \code{fullMerge}. The core definitions do not need
to change for that runner.

\subsection{Commit Graphs}

\code{PushHistory} is a linear list of \code{PatchStep}s. It does not model
branches or merge commits. A graph model needs a commit table, parent commit
ids, and a nearest-common-ancestor function.

\subsection{Merged Diff Metadata}

\code{Diff} stores one \code{agent} and one \code{timestamp}. A clean merge may
combine edits from multiple agents. A later version can add a field that lists
all contributing agents and timestamps.

\subsection{Timestamp Parser}

Lean stores \code{timestamp} as \code{String}. A later version can add a
\code{Timestamp} type and reject strings that are not valid ISO 8601
date-times.

\subsection{Sequence Numbers}

\code{seq} is checked to be positive. It is not checked to be exactly
\code{1, 2, ..., n}. A later schema check can add that rule.

\subsection{File IO}

The core code has no file-system dependency. A separate executable can read a
file, collect diff JSON, validate it, and print either a merge result or
conflict witnesses.
